% =====================================================================
%  CV Harshil Gorasiya (English)
%  Kompilieren mit pdflatex (zweimal). Legen Sie photo.png in denselben Ordner.
%  Datei fuer den Versand umbenennen: Lebenslauf_Harshil_Gorasiya.pdf
%  Layout bewusst ATS-freundlich: Standardschrift, Klartext, keine Icons,
%  keine Fortschrittsbalken, einspaltiger Textfluss, eindeutige Ueberschriften.
% =====================================================================
\documentclass[10pt,a4paper]{article}

\usepackage[a4paper,left=1.7cm,right=1.7cm,top=1.1cm,bottom=1.35cm]{geometry}
\usepackage[T1]{fontenc}
\usepackage[utf8]{inputenc}
\usepackage{textcomp}
\usepackage[scaled=0.96]{helvet}
\renewcommand{\familydefault}{\sfdefault}
\usepackage[english]{babel}
\usepackage{microtype}
\usepackage{xcolor}
\usepackage{graphicx}
\usepackage{enumitem}
\usepackage{titlesec}
\usepackage{array}
\usepackage{fancyhdr}
\usepackage{lastpage}
\makeatletter

\definecolor{navy}{HTML}{1B365D}
\definecolor{mid}{HTML}{5A6472}
\definecolor{tint}{HTML}{EEF2F7}

\usepackage[unicode,colorlinks=true,urlcolor=navy,linkcolor=navy,
  pdftitle={CV Harshil Gorasiya},
  pdfauthor={Harshil Gorasiya},
  pdfsubject={CV: AI/ML Engineering, LLM, RAG, AI agents},
  pdfkeywords={working student, master thesis, AI, machine learning, LLM, RAG, Python, PyTorch}]{hyperref}

\input{glyphtounicode}
\pdfgentounicode=1

% Layout bleibt in jeder Umgebung identisch (keine Silbentrennung, Flattersatz)
\hyphenpenalty=10000
\exhyphenpenalty=10000
\emergencystretch=3em
\setlength{\parindent}{0pt}
\setlength{\parskip}{0pt}
\raggedright
\newlength{\leftcol}\setlength{\leftcol}{2.45cm}
\newcommand{\gap}{0.5em}
\setlist[itemize]{leftmargin=*,labelindent=\dimexpr\leftcol+\gap\relax,labelwidth=0.55em,labelsep=0.4em,align=left,itemsep=0.6pt,topsep=1.5pt,parsep=0pt,partopsep=0pt,label={\color{navy}\textbullet}}


% Fusszeile
\pagestyle{fancy}
\fancyhf{}
\renewcommand{\headrulewidth}{0pt}
\renewcommand{\footrulewidth}{0pt}
\fancyfoot[C]{\footnotesize\color{mid}Harshil Gorasiya \textperiodcentered{} CV \textperiodcentered{} Page \thepage{} of \pageref{LastPage}}

% Abschnittsueberschriften
\titleformat{\section}{\normalfont\bfseries\color{navy}\normalsize}{}{0pt}
  {\MakeUppercase}[\vspace{-0.45em}{\color{navy}\rule{\linewidth}{0.7pt}}]
\titlespacing*{\section}{0pt}{11pt}{5pt}

% Hilfsmakros
\newcommand{\sep}{\hspace{0.45em}\textbar\hspace{0.45em}}
\newcommand{\dotsep}{\hspace{0.3em}\textperiodcentered\hspace{0.3em}}
\newcommand{\entry}[3]{%
  \begingroup\raggedright\leftskip=\dimexpr\leftcol+\gap\relax\noindent
  \hspace*{-\dimexpr\leftcol+\gap\relax}%
  \vtop to 0pt{\hsize=\dimexpr\leftcol+\gap\relax\raggedright\small\color{mid}\noindent #1\par\vss}%
  #2\par #3\par\endgroup\vspace{5pt}}
\newcommand{\etitle}[2]{\textbf{#1}\ifx&#2&\else\textcolor{mid}{\sep}#2\fi}
\newcommand{\stack}[1]{{\small\itshape\color{mid}#1}\par\vspace{5pt}}
\newcommand{\kpi}[2]{\begin{minipage}[t]{\linewidth}\centering{\large\bfseries\color{navy}#1}\\[-1pt]{\scriptsize\color{mid}#2}\end{minipage}}

\renewcommand\normalsize{\@setfontsize\normalsize{10.4}{12.6}}
\makeatother
\begin{document}
\normalsize

\noindent
\begin{minipage}[t]{\dimexpr\textwidth-3.05cm\relax}
  {\fontsize{25}{28}\selectfont\bfseries\color{navy}Harshil Gorasiya}\\[4pt]
  {\bfseries LLM systems: retrieval, AI agents, and small-model research}\\[1pt]
  {\color{navy}\bfseries Python\dotsep PyTorch\dotsep LLM\dotsep RAG\dotsep Computer Vision}\\[5pt]
  {\small
  Heilbronn, Germany\sep +49 155 6351 7346\sep
  \href{mailto:harshil.gorasiya.0011@gmail.com}{harshil.gorasiya.0011@gmail.com}\\[1pt]
  LinkedIn: \href{https://www.linkedin.com/in/harshil-gorasiya-054a58319}{linkedin.com/in/harshil-gorasiya-054a58319}\\[1pt]
  GitHub: \href{https://github.com/N0t-Harshil}{github.com/N0t-Harshil}\sep
  Portfolio: \href{https://harshil-0011.github.io/My-Portfolio/}{harshil-0011.github.io/My-Portfolio}}
\end{minipage}\hfill
\begin{minipage}[t]{2.7cm}
  \raggedleft
  \IfFileExists{photo.png}{\setlength{\fboxsep}{0pt}\setlength{\fboxrule}{0.6pt}%
    \fcolorbox{navy}{white}{\includegraphics[width=2.7cm]{photo.png}}}{}
\end{minipage}

\vspace{7pt}
\noindent\colorbox{tint}{\parbox{\dimexpr\textwidth-2\fboxsep\relax}{\small
  \textbf{Seeking:} Working student position (20 hrs/week during semester)\dotsep
  Full-time during semester breaks\dotsep Entry-level LLM/RAG role from 2027\\[1pt]
  \textbf{Available:} immediately\sep
  \textbf{Residence status:} valid residence permit with work permit}}

\vspace{6pt}
\noindent
\begin{tabular*}{\textwidth}{@{\extracolsep{\fill}}*{5}{>{\centering\arraybackslash}p{0.185\textwidth}}}
  \kpi{+8 pts}{accuracy: 71\,\% to 79\,\%} &
  \kpi{–20\,\%}{training time (V100 GPUs)} &
  \kpi{2.4\texttimes}{throughput vs. YOLOv5s} &
  \kpi{28--32 FPS}{real-time object detection on CPU} &
  \kpi{8 providers}{LLMs with failover (Ardan-CLI)} \\
\end{tabular*}

\section*{Profile}
Master's student at Hochschule Heilbronn. I build LLM applications: local RAG systems with source citations, AI agents with tool use, and computer vision models that run on CPU. I also bring two ML internships, 10 months of AI-enabled web development, and a research project on reasoning modules for language models.

\section*{Education}
\entry{since 03/2025}
  {\etitle{M.Sc. Software Engineering and Management}{Hochschule Heilbronn}}
  {\small English-taught programme\sep GPA so far: 1.9 (German scale)\sep expected graduation: 2027}
\entry{2020 -- 2024}
  {\etitle{B.Tech. Information Technology}{A.\,D. Patel Institute of Technology, Gujarat}}
  {\small India\sep Grade: 1.9 (converted to German scale)}

\section*{Projects}
\entry{AI agent}
  {\etitle{Ardan-CLI: Autonomous Terminal Coding Agent}{}}
  {\stack{TypeScript, Node.js, MCP, LSP}
  \begin{itemize}
    \item Agent with streaming, plan mode, sub-agents, and checkpoints; \textbf{35 tools}.
    \item \textbf{8 LLM providers} with automatic failover; \textbf{0 runtime dependencies}, MCP and LSP clients built by hand.
    \item \textbf{8.1 ms} overhead per agent step (model time excluded); automated tests on Linux, macOS, and Windows.
  \end{itemize}}
\entry{Research engine}
  {\etitle{Local Perplex: Local AI Research Assistant with Citations}{}}
  {\stack{Python, FastAPI, SSE, SQLite, Ollama, Docker}
  \begin{itemize}
    \item Searches the web and your own documents (PDF, Word, Excel) and streams answers with citations.
    \item \textbf{26 LLM providers}, 7 search backends, fully offline if desired.
    \item \textbf{17\texttimes{} better retrieval quality} than my first version (BEIR SciFact); answer page in about \textbf{200 ms}; knowledge graph with \textbf{77\texttimes{}} faster writes.
  \end{itemize}}
\entry{RAG system}
  {\etitle{ANVITHA: Local Knowledge Vault with Grounded, Cited Answers}{}}
  {\stack{Python, FastAPI, Ollama, SQLite, BM25}
  \begin{itemize}
    \item Turns PDF, text, and Markdown files into a linked knowledge vault; every answer cites the page of the original source.
    \item Each answer sentence is checked against the sources; unsupported claims are removed.
    \item Questions outside the topic are refused in under \textbf{0.3 s}, without a model call.
    \item Notes stored as editable Markdown with a full change log; runs fully local via Ollama.
  \end{itemize}}
\newpage
\entry{Master's thesis\\(ongoing)}
  {\etitle{Ekdant: Reasoning Module for Language Models}{}}
  {\stack{Hochschule Heilbronn, supervised by a professor\sep PyTorch, Hugging Face}
  \begin{itemize}
    \item Parallel debate module (4 paths, consensus) added as a final stage to SmolLM2-135M; \textbf{+19.6M parameters}.
    \item Controlled comparison against an architecture-matched baseline on medical question answering: the best variant leads on every metric, by a narrow margin.
    \item Literature base of 66 research papers.
  \end{itemize}}
\entry{More projects}
  {\etitle{AI Video Upscaling}{Python, PyTorch}}
  {\small Windows desktop application for video super-resolution with model selection; compared Real-ESRGAN variants using PSNR and SSIM.}
\entry{\strut}
  {\etitle{MNIST-MLOps}{Python, PyTorch, Flask, SQLite}}
  {\small Incremental training, model version management, and a web dashboard.}

\section*{Professional Experience}
\entry{05/2024 -- 02/2025\\(10 months)}
  {\etitle{AI \& Web Solutions Developer}{Aani Brothers Infotech, Surat, India}}
  {\begin{itemize}
    \item Built web applications that turn business requirements into user interfaces and workflows, including AI-enabled features.
    \item Conducted user research, wireframing, and prototyping; optimized designs for usability.
  \end{itemize}}
\entry{12/2023 -- 04/2024\\(5 months)}
  {\etitle{Deep Learning Intern}{PD Avenue, Surat, India}}
  {\begin{itemize}
    \item Built a residual-block CNN (TensorFlow/Keras) for age and gender classification on \textbf{22,000+ images}; raised accuracy from \textbf{71\,\% to 79\,\%}.
    \item Cut training time on V100 GPUs by \textbf{20\,\%} (learning-rate scheduling, larger batch size) with no loss of accuracy.
    \item Identified and fixed a data-leakage issue.
  \end{itemize}}
\entry{05/2023 -- 06/2023\\(2 months)}
  {\etitle{Machine Learning Intern}{Dotcom IoT LLP, Surat, India}}
  {\begin{itemize}
    \item YOLOv8n object detection on a CPU-only edge device (Intel Core i7): 720p video at \textbf{28--32 FPS}; \textbf{2.4\texttimes{} higher throughput} than YOLOv5s.
    \item Curated and annotated \textbf{3,000+ images} across five vehicle classes.
    \item Reduced false positives in low light by \textbf{18\,\%} through IoU and NMS tuning.
  \end{itemize}}

\section*{Skills}
\entry{Programming}
  {\small Python, TypeScript, Node.js, SQL, Bash}{}
\entry{Engineering}
  {\small Git, GitHub Actions (CI/CD), Docker, Linux, pytest, REST APIs, FastAPI, Flask, streaming (SSE), JSON-RPC}{}
\entry{ML \& vision}
  {\small PyTorch, TensorFlow/Keras, Hugging Face Transformers, YOLOv8, CNNs with residual blocks, object detection, super-resolution (Real-ESRGAN), OpenCV, data annotation, hyperparameter tuning, learning-rate scheduling, overfitting analysis}{}
\entry{LLM \& agents}
  {\small RAG systems, retrieval (BM25, embeddings), knowledge graphs, source-grounded answering, prompt engineering, fine-tuning, small language models (SmolLM2), AI agents with tool use (ReAct), Model Context Protocol (MCP), Anthropic API, OpenAI-compatible APIs, Ollama}{}
\entry{Evaluation \& research}
  {\small A/B experiment design, model evaluation (accuracy, latency, throughput, PSNR/SSIM, perplexity, ROUGE-L), literature review}{}
\entry{Databases \& UX}
  {\small SQLite, PostgreSQL, MongoDB, NumPy\sep user research, wireframing, prototyping, usability}{}

\section*{Certifications}
\entry{2026}{\small Claude with the Anthropic API\sep Model Context Protocol: Advanced Topics (both Anthropic)}{}
\entry{2025}{\small Machine Learning with Python (IBM)\sep Fundamentals of Data Science (Databricks)}{}

\section*{Languages}
\entry{\strut}{\small \textbf{English:} fluent\sep \textbf{German:} A2--B1 (self-assessed)\sep \textbf{Hindi:} fluent\sep \textbf{Gujarati:} native}{}

\vspace{2pt}
{\small\color{mid}Heilbronn, 05.11.2026}

\end{document}
